% /*@@
%   @file      Preface.tex
%   @date      27 Jan 1999
%   @author    Tom Goodale, Gabrielle Allen, Gerd Lanferman
%   @desc 
%   Preface for the Cactus User's Guide
%   @enddesc 
%   @version $Header$      
%   @history
%   @date       Sun Jul 20 12:57:28 CEST 2003
%   @author     Jonathan Thornburg
%   @desc       split off reference material into a separate
%               Reference Manual (in new directory ../ReferenceManual/)
%   @endhistory
% @@*/

%%%%%%%%%%%%%%%%%%%%%%%%%%%%%%%%%%%%%%%%%%%%%%%%%%%%%%%%%%%%%%%%%%%%%%%%%%%%%%%%
%%%%%%%%%%%%%%%%%%%%%%%%%%%%%%%%%%%%%%%%%%%%%%%%%%%%%%%%%%%%%%%%%%%%%%%%%%%%%%%%

{\large \bf Preface} 
\label{sec:pr}
 
\vskip .5cm

This document will eventually be a complete guide to using and
developing with the Cactus Code. However, it is currently under
development, so please be patient if you can't find what you need.
Please report omissions, errors, or suggestions to 
and of our contact addresses below, and we will
try and fix them as soon as possible. 

\vskip .5cm

\textbf{Overview of documentation}

\vskip .5cm

This guide covers the following topics

\begin{Lentry}

\item [\textbf{Part~\ref{part:RunningCactus}: Installation and Running.}]
  A description of required hardware and software and an installation guide
  for Cactus. 
  Describes how to check the installation with Cactus test suites. 

\item [\textbf{Part~\ref{part:ThornWriters}: Application Thorn Writing.}]
  An introduction to thorn concepts and description of how to create, write
  and maintain application thorns. Explanation of use of the
  programming interface to take advantage of parallelism and modularity.

\item [\textbf{Part~\ref{part:UtilityRoutines}: Utility Routines.}]
  A discussion of low-level utility routines available in Cactus, such as
  key/value tables.

\item [\textbf{Part~\ref{part:Infrastructure}:
       Infrastructure Thorn Writer's Guide.}]
  A more advanced discussion of user supplied infrastructure routines such
  as additional \textit{output routines}, \textit{drivers}, etc.
  [\emph{This part is currently under construction.}]

\item [\textbf{Part~\ref{part:Appendices}: Appendices.}]
        These contain a description of the Cactus Configuration Language,
        a glossary,
        and other odds and ends, such as how to use GNATS and TAGS.

\end{Lentry}

Related topics are discussed in separate documents including:

\begin{Lentry}

\item [\textbf{Reference Manual}] Contains detailed descriptions of
  the functions provided by the Cactus flesh API, along with
  other reference material.

%\item [{\bf Computational Thorn Guide}] This will contain details about the 
%arrangements and thorns making up the standard Cactus Computational Tool Kit.

%\item [{\bf Relativity Thorn Guide}] This will contain details about the arrangements and thorns making up the Cactus Relativity Tool Kit, one of the major 
% motivators, and still the driving force, for the Cactus Code.

%\item [{\bf Flesh Maintainers Guide}] 
% This will contain all the gruesome details
% about the inner workings of Cactus, for all those who want or need to 
% expand or maintain the core of Cactus.

\end{Lentry}

\vskip .5cm

{\bf Typographical Conventions}

\begin{Lentry}

\item[{\tt Typewriter}] Is currently used for everything you type,
	for program names, and code extracts.
\item[{\tt < ... >}] Indicates a compulsory argument.
\item[{\tt [ ... ]}] Indicates an optional argument.
\item[{\tt |}] Indicates an exclusive or.

\end{Lentry}
 
\vskip .5cm

{\bf How to Contact Us}

\vskip .5cm

Please let us know of any errors or omissions in this guide, as well
as suggestions for future editions. These can be reported via our 
bug tracking system at \url{http://www.cactuscode.org}, or via email to
\code{cactusmaint@cactuscode.org}. Alternatively, you can write to us at

\vskip .5cm
The Cactus Team\\
Center for Computation \& Technology\\
216 Johnston Hall\\
Louisiana State University\\
Baton Rouge, LA 70803\\
USA


\vskip .5cm

{\bf Acknowledgements}

\vskip .5cm

Hearty thanks to all those who have helped with documentation for the
Cactus Code. Special thanks to those who struggled with the earliest
sparse versions of this guide and sent in mistakes and suggestions,
in particular John Baker, Carsten Gundlach, Ginny Hudak-David, 
Sai Iyer, Paul Lamping, Nancy Tran and Ed Seidel. 

%%%%%%%%%%%%%%%%%%%%%%%%%%%%%%%%%%%%%%%%%%%%%%%%%%%%%%%%%%%%%%%%%%%%%%%%%%%%%%%%
%%%%%%%%%%%%%%%%%%%%%%%%%%%%%%%%%%%%%%%%%%%%%%%%%%%%%%%%%%%%%%%%%%%%%%%%%%%%%%%%
